% Introduccion - la escribe D al integrar (media pagina como maximo)

\section{Introducción}

El taller pide tres cosas sobre el mismo sistema del Proyecto 1: prototipar en
alto nivel las frecuencias con las que se va a trabajar, generar en el
microcontrolador el emisor de esas señales, y definir con diagramas cómo se va a
validar lo que el receptor recibe. No es un ejercicio aparte: lo que aquí se
decide es el banco de pruebas con el que después se va a medir el sistema
distribuido de reducción de ruido.

El montaje tiene dos placas. Una Raspberry Pi Pico hace de emisor y sintetiza por
PWM, con un filtro RC a la salida, tres modos de señal; una Pico W hace de nodo 1
y las captura por su ADC, las transforma por bloques y reporta lo que midió. Los
tres modos no se eligieron al azar: el modo 1 es un tono de 440\,Hz, el modo 2 el
multitono de 440, 880 y 1320\,Hz, y el modo 3 ese multitono más una interferencia
de 3500\,Hz con la misma relación señal a ruido del \texttt{caso\_1} de la Tarea 1.
Así cada modo corresponde a una prueba que el proyecto ya tenía definida, y los
resultados del taller se pueden comparar directamente con los que se midieron en
simulación.

Ningún número de este documento está escrito a mano en el código. Los parámetros
que vienen de la Tarea 1 --- frecuencia de muestreo de 16\,kHz, bloques de 1024
muestras, ventana Hann y 50\,\% de solapamiento --- se leen de
\texttt{metadatos/configu\-raciones.json}, y los que agrega el taller --- los modos
del emisor, los datos del receptor y los márgenes de validación --- de
\texttt{taller\_semana12/parametros.json}. Si en la defensa se cambia un
parámetro, se cambia en el JSON y todo lo demás lo sigue.

El documento está organizado por bloques, en el orden en que se recorre la señal.
La sección \ref{sec:senales} justifica la frecuencia de muestreo y el tamaño de
bloque y define la señal exacta de cada modo, que es la referencia que usan todos
los demás. La sección \ref{sec:sistema} arma el sistema completo en alto nivel. La
sección \ref{sec:emisor} lo baja al microcontrolador. La sección
\ref{sec:validacion} define la estrategia de validación en el receptor con sus
cuatro diagramas y los criterios de aceptación. La sección \ref{sec:discusion}
cierra con lo que funcionó, las limitaciones y lo que cambia para el proyecto.
